\documentclass{situationist-preprint}

\title{Responsibility for Omissions}
\msid{ACT-07-a}
\msfamily{ACT-07}
\msdiscipline{Action, Agency and Free Will}
\mstype{Principal manuscript}
\mskeywords{omissions; moral responsibility; ability; obligation; Frankfurt cases; Sharks case; asymmetry}

\begin{document}
\maketitle

\begin{abstract}
We are blamed for what we fail to do as well as for what we do: for not rescuing the drowning child, for not paying the bill, for not warning the friend. But omissions are absences, and absences are many; it is unclear what makes some of them, and not others, things for which an agent answers. We argue that an agent is responsible for an omission when she could and should have acted: when she had the ability to perform the omitted action and was under a requirement, moral or other, to perform it, and when her not performing it issues from her own reasons-responsive agency. The key move is to distinguish an omission, which is a failure to exercise an ability one has, from a mere non-performance. It follows that Frankfurt cases, which show that responsibility for actions does not require alternative possibilities, do not carry over to omissions: the devices that remove the alternative also remove the ability whose non-exercise the omission consists in. An agent in such a case remains responsible for her decision not to act, but not for the outcome she could not have prevented.
\end{abstract}

\section{The problem}

A man walks along a beach and sees a child drowning. He could easily save her, judges that he should, and decides not to bother. He is blameworthy for not saving her. Ordinary moral practice treats such failures as on a par with actions: the law recognizes liability for omissions where there is a duty to act, and morality recognizes it more widely. Two questions arise. The first is which non-performances are omissions for which an agent can be responsible, given that at any moment each agent fails to do indefinitely many things. The second is whether the conditions of responsibility for omissions are those for actions.

The second question has been sharpened by Frankfurt cases. \citet{Frankfurt1969} argued that an agent may be responsible for an action although she could not have done otherwise, since a counterfactual intervener, who never in fact intervenes, would have ensured that she acted as she did. ACT-02 defends this conclusion for actions. \citet{vanInwagen1978} argued that analogues of the principle of alternative possibilities hold for omissions and for outcomes even if they fail for actions. \citet{FischerRavizza1991} gave the case that has framed the debate since. John, walking along the beach, decides not to save a drowning child; unknown to him, the water is full of sharks that would have prevented him from reaching her. John is not responsible, they argued, for not saving the child, though he is responsible for not trying and for his decision. \citet{Frankfurt1994} replied that the asymmetry is only apparent, and that cases of the Frankfurt type can be built for omissions as for actions. \citet{Sartorio2005} argued for a different asymmetry between actions and omissions, and \citet{Clarke2014} has given the most complete treatment of the metaphysics of omissions and of responsibility for them. We side with Fischer and Ravizza on the Sharks case, and give a principled account of why.

\section{Framework}

\begin{definition}[Non-performance and omission]\label{d:omission}
An agent \emph{does not perform} $A$ at $t$ if she does not do $A$ at $t$. She \emph{omits} $A$ at $t$ if she does not perform $A$ at $t$, she has at $t$ the ability to perform $A$, and there is at $t$ a requirement or standard applying to her that calls for $A$.
\end{definition}

The definition follows the ordinary use of \enquote{omission}, on which one omits only what one could have done and was in some way expected to do; it is the narrow sense that \citet{Clarke2014} distinguishes from the wider sense of mere not-doing. Indefinitely many non-performances are not omissions. The man on the beach does not perform a sonata; he does not thereby omit to perform one.

\begin{definition}[Ability]\label{d:ability}
An agent has at $t$ the \emph{ability} to perform $A$ if she has the general capacity to do $A$, and the opportunity: nothing in her circumstances at $t$ would prevent her from doing $A$ if she tried.
\end{definition}

\begin{principle}[Should]\label{p:should}
An agent can be blameworthy for not performing $A$ only if she was under a requirement to perform $A$, and only if she was in a position to know, or culpably failed to be in a position to know, that she was.
\end{principle}

\begin{principle}[Ownership]\label{p:own}
An agent is responsible for what she does or fails to do only if it issues from a mechanism that is her own and reasons-responsive, in the sense of ACT-01-a.
\end{principle}

Principle~\ref{p:should} states the normative condition: a non-performance can make one blameworthy only if something was called for. The requirement may be a duty of rescue, which ETH-18 argues we have up to moderate cost, a role obligation, a promise, or a legal duty. The epistemic clause is the one ETH-23 states for past agents: agents are blameworthy only for what they could reasonably have known \citep{Smith1983}. Principle~\ref{p:own} is the control condition that ACT-01-a defends for free action.

\section{The argument}

\begin{lemma}\label{l:could}
An agent is responsible for an omission only if she could have performed the omitted action.
\end{lemma}

\begin{proof}
To be responsible for an omission is to be responsible for one's failure to do $A$. By Definition~\ref{d:omission}, an agent omits $A$ only if she had the ability to perform $A$; where she lacked it, there is a non-performance of $A$ but no omission of $A$, and so nothing of that kind for her to be responsible for. The failure for which she answers is a failure to exercise an ability she had; an agent who lacked the ability did not fail to exercise it.
\end{proof}

Lemma~\ref{l:could} is not an application of the principle of alternative possibilities to omissions. It does not say that an agent is responsible for something only if she could have done otherwise than it. It says that the object of responsibility, in the case of an omission, is constituted in part by an ability, so that where the ability is absent the object is absent.

\begin{theorem}\label{t:main}
An agent is responsible for omitting $A$ if, and only if, she could have performed $A$, she should have performed $A$, and her non-performance of $A$ issued from her own reasons-responsive mechanism.
\end{theorem}

\begin{proof}
Necessity: the first condition is Lemma~\ref{l:could}; the second is Principle~\ref{p:should}, applied to responsibility in the sense that makes one an apt target of blame; the third is Principle~\ref{p:own}. Sufficiency: suppose all three hold. The agent then omitted $A$ in the sense of Definition~\ref{d:omission}. Her non-performance issued from her own practical reasoning, which recognized or could have recognized the requirement and could have issued in her doing $A$ in a range of situations; this is what it is for her failure to be under her control in the sense ACT-01-a defends. An agent who fails, under her own control, to do what she had the ability to do and was required to do, knowing or in a position to know of the requirement, satisfies every condition of responsibility that she would satisfy for an action done under her own control. So she is responsible for the omission.
\end{proof}

\begin{corollary}\label{c:frankfurt}
Frankfurt cases do not carry over to omissions. In a case in which a device or circumstance would have prevented the agent from performing $A$, she is not responsible for omitting $A$, though she may be responsible for her decision not to perform $A$ and for not trying.
\end{corollary}

\begin{proof}
A Frankfurt-style case for omissions would be one in which the agent does not do $A$, could not have done $A$, and is nonetheless responsible for not doing $A$. In the Sharks case, the sharks remove John's opportunity, so that he lacks the ability to save the child (Definition~\ref{d:ability}). By Lemma~\ref{l:could}, he is not responsible for omitting to save her. He did have the ability to decide to save her and to try: nothing prevented him from jumping in. These are actions, or omissions of actions he could perform, and he is responsible for deciding as he did and for not trying, by Theorem~\ref{t:main} for the latter and, for the former, by the account of responsibility for actions that ACT-02-a defends.
\end{proof}

The asymmetry with actions is now explained. In a Frankfurt case for actions, the intervener removes the agent's alternatives, but the action she performs is still hers: it is constituted by her deliberation, decision and movement, none of which depend on the alternatives. ACT-02-a states this as the principle that responsibility for an action depends only on its actual sequence. In the omission case, what the device removes is the ability whose non-exercise the omission is. The \enquote{actual sequence} of an omission includes the agent's standing ability to act, because an omission is a failure of that ability to be exercised.

\section{Objections and replies}

\begin{objection}[Frankfurt's symmetry]
\citet{Frankfurt1994} held that the Sharks case is mis-described: John is responsible for not saving the child in the sense that matters, since his not saving her expresses exactly the indifference for which he is blamed, and the sharks play no role in what happens.
\end{objection}

\begin{reply}
We agree that the sharks play no role in John's deliberation, and that he is blameworthy for the indifference his decision expresses. Corollary~\ref{c:frankfurt} preserves all of this: John is responsible for his decision and for not trying. What it denies is that he is responsible for the outcome, the child's not being saved, and that he should be blamed as one who could have saved her. Consider the person who would have saved the child if she could, and was prevented by the sharks: she is not blamed at all, and John is blamed more only for what he could control, his decision. Were he responsible for not saving the child, his blame would be that of the man on a shark-free beach, who could have saved her; it is not.
\end{reply}

\begin{objection}[Overgeneration]
The account makes agents responsible for too many omissions. Every affluent person could give to a famine fund and arguably should; Theorem~\ref{t:main} makes each responsible for countless omissions every day.
\end{objection}

\begin{reply}
The theorem makes an agent responsible for an omission only where she was under a requirement to act, and its extent is fixed by first-order morality, not by the theory of responsibility. If the duties of beneficence are as ETH-06 argues, demanding a substantial fixed share of income and no more, then an agent who has given that share omits nothing in not giving more, since nothing further was required of her. Where there is a duty and she fails, it is the right result that she is responsible.
\end{reply}

\begin{objection}[Causal idleness]
Omissions are absences, and absences cause nothing. If responsibility for an outcome requires causing it, no one is responsible for the outcomes of omissions, such as the child's drowning.
\end{objection}

\begin{reply}
Theorem~\ref{t:main} concerns responsibility for the omission, not for its outcome. Responsibility for an outcome of an omission requires in addition that the outcome depend on it: that had the agent done $A$, the outcome would not have occurred. This is a counterfactual, not a productive, relation, and whether it suffices for causation is a question about causation that we need not settle. Ordinary judgments that the negligent lifeguard is responsible for the drowning track the counterfactual dependence together with the conditions of the theorem.
\end{reply}

\begin{objection}[Unwitting omissions]
The person who forgets her friend's birthday did not decide anything; no mechanism of practical reasoning issued in her failure. Yet she may be blameworthy.
\end{objection}

\begin{reply}
Her non-performance issued from a mechanism that is hers and that would have responded to reasons she had: had she attended to her calendar, as she could and should have, she would have remembered. Principle~\ref{p:own} does not require a decision, only that the failure be under the control of a reasons-responsive mechanism. Where the forgetting is not of this kind, as when it results from illness, the agent is not responsible, and that is the right verdict.
\end{reply}

\section{Conclusion}

An omission is a failure to exercise an ability one has, when something calls for its exercise. An agent is responsible for an omission when she could and should have acted, and her failure issues from her own reasons-responsive agency. Because the ability to act belongs to what an omission is, Frankfurt cases, which preserve the action while removing alternatives, do not carry over: the device that removes the alternative removes the omission with it. We can be responsible for omissions, under the same control condition as for actions, but with an ability condition that actions do not have.

\begin{thebibliography}{99}\small
\bibitem[Clarke(2014)]{Clarke2014} Clarke, R. (2014). \emph{Omissions: Agency, Metaphysics, and Responsibility}. Oxford University Press.
\bibitem[Fischer and Ravizza(1991)]{FischerRavizza1991} Fischer, J.~M. and Ravizza, M. (1991). Responsibility and inevitability. \emph{Ethics}, 101(2), 258--278.
\bibitem[Frankfurt(1969)]{Frankfurt1969} Frankfurt, H.~G. (1969). Alternate possibilities and moral responsibility. \emph{The Journal of Philosophy}, 66(23), 829--839.
\bibitem[Frankfurt(1994)]{Frankfurt1994} Frankfurt, H.~G. (1994). An alleged asymmetry between actions and omissions. \emph{Ethics}, 104(3), 620--623.
\bibitem[Sartorio(2005)]{Sartorio2005} Sartorio, C. (2005). A new asymmetry between actions and omissions. \emph{No\^us}, 39(3), 460--482.
\bibitem[Smith(1983)]{Smith1983} Smith, H. (1983). Culpable ignorance. \emph{The Philosophical Review}, 92(4), 543--571.
\bibitem[van Inwagen(1978)]{vanInwagen1978} van Inwagen, P. (1978). Ability and responsibility. \emph{The Philosophical Review}, 87(2), 201--224.
\end{thebibliography}

\end{document}
